% Source: source/普通高考/2012/2012江苏.pdf, page 1, question 4.
% Type: algorithm flowchart; pdfimages -list reports no embedded bitmap (vector line art).
% Grid evidence: tmp/review_flowchart_2012_jiangsu/grid/page-1-grid100.png
% (100px coarse + 50px fine) and q4-block-grid50.png (50px coarse + 25px fine).
% Comparison crop: tmp/review_flowchart_2012_jiangsu/crop/q4_original_final.png,
% taken from the ungridded page render with a safety border.
% Node-edge list: 开始→k←1→(k^2−5k+4>0)；是→输出k→结束；
% 否→k←k+1→回线左向箭头→判定前共享主干。All borders/connectors are solid;
% labels Y/N are placed beside the corresponding exits. Nodes are centered with
% local parindent=0pt; the return arrow and the main downward arrow are distinct.
\begin{tikzpicture}[
  x=1cm,
  y=0.9cm,
  >=Stealth,
  line width=0.45pt,
  line cap=round,
  line join=round,
  every node/.style={font=\normalsize, inner sep=1pt, align=center, anchor=center,
    execute at begin node={\setlength{\parindent}{0pt}\noindent}},
  flow/.style={draw, line width=0.45pt},
  terminal/.style={flow, rounded corners=3pt, minimum width=1.25cm, minimum height=0.70cm},
  process/.style={flow, minimum height=0.72cm, align=center},
  io/.style={flow, trapezium, trapezium left angle=74, trapezium right angle=106,
    minimum height=0.76cm, align=center},
  decision/.style={flow, diamond, aspect=2.9, minimum width=4.0cm,
    minimum height=1.0cm, inner sep=1pt, align=center}
]
  \node[terminal] (start) at (0,8.20) {开始};
  \node[process, minimum width=1.50cm] (init) at (0,6.80) {$k\gets1$};
  \node[decision] (test) at (0,4.95) {$k^2-5k+4>0$};
  \node[io, minimum width=2.05cm] (output) at (0,3.35) {输出 $k$};
  \node[terminal] (finish) at (0,1.95) {结束};
  \node[process, minimum width=2.35cm] (increment) at (3.60,4.95) {$k\gets k+1$};

  \coordinate (loopJoin) at (0,6.10);
  \coordinate (loopRight) at (3.60,6.10);
  \coordinate (loopTip) at (0.35,6.10);

  % Main path and terminating branch.
  \draw[->] (start.south) -- (init.north);
  \draw (init.south) -- (loopJoin);
  \draw[->] (loopJoin) -- (test.north);
  \draw[->] (test.south) -- node[right=2pt] {Y} (output.north);
  \draw[->] (output.south) -- (finish.north);

  % N branch and return loop. The left-pointing return arrow ends on the
  % shared stem; the main stem retains its single downward arrow into test.
  \draw[->] (test.east) -- node[above=3pt, pos=0.18] {N} (increment.west);
  \draw[->] (increment.north) -- (loopRight) -- (loopTip) -- (loopJoin);
  \path[use as bounding box] (-2.25,1.55) rectangle (4.95,8.60);
\end{tikzpicture}
